FANTOM’s performance deteriorates when a regime contains only a handful of samples or is recorded at an extremely low sampling rate. This shortcoming is not surprising, estimating a separate causal graph for each regime is intrinsically difficult in the presence of multiple regimes and heteroscedastic noise. Yet this ability to pinpoint which edges vanish or emerge from one regime to the next is what makes FANTOM valuable in domains such as healthcare and climate science, where regime shifts carry substantive meaning. Importantly, in realistic settings where each regime offers sufficient data, for example, epileptic seizures that last several minutes at 250 Hz, FANTOM delivers strong  results and provides insights unattainable with existing methods.

FANTOM requires a suitable initial segmentation to effectively learn the regime indices. Our ablation studies demonstrate that selecting a reasonable initial window size establishes a basis for accurate regime detection, which is crucial for achieving high performance. However, we highlight a critical limitation regarding the subsequent pruning step: if the initial window is too high (over-pruning), it can drastically reduce the number of initial regimes below the ground truth. This initialization leads to a loss of essential information and a significant deterioration of FANTOM’s final performance.

Spurious causality is a practical risk, so we advocate using the learned graphs as decision support rather than as ground truth. Before acting on them, especially in medical or financial settings, practitioners should adopt safeguards: (i) expert vetting of proposed edges and mechanisms; (ii) robustness checks (e.g., stability under resampling or perturbations, alternative specifications); and (iii) interventional validation when feasible.